\documentclass[a4paper,12pt]{article}

\usepackage{sbc-template}
\usepackage{graphicx,url}
\usepackage[utf8]{inputenc}
\usepackage{fancyhdr}
\usepackage[brazil]{babel}
\usepackage{hyperref}
\usepackage{xcolor}
\usepackage{nameref}
%\usepackage[latin1]{inputenc}  

\pagenumbering{gobble}                          % sets page numbering off
\pagestyle{fancy}\setlength\headheight{100pt}   % set the pagestyle to fancy to include headers and footers
\pagestyle{fancy}\setlength\textheight{650pt}
\pagestyle{fancy}\setlength\footskip{20pt}
\fancyhf{}                                      % clear the header and footer
\renewcommand{\headrulewidth}{0pt}              % no horizontale ruler under header
\fancyhead[L]{
\includegraphics[width=3.8cm]{img/LogoINCT.jpg}}
\fancyhead[R]{
\includegraphics[width=3.8cm]{img/LogoOdissea.jpg}}

\begin{document} 

%\title{Guia para o carregamento de dados no Geonode Odisseia}
%\date{}
%\maketitle
%\thispagestyle{fancy}

\begin{center}
  \vspace{12pt}
  \Large\textbf{Guia para construção de Dashboards no Geonode}
\end{center}

\section{Introdução}

  Este documento explicita as funcionalidades dos Dashboards do Geonode e traz
  instruções para o seu uso. As informações aqui contidas são provenientes das
  documentações oficiais do 
  \textcolor{blue}
  {\href{https://docs.geonode.org/en/master/usage/dashboard/index.html}
  {Geonode}}
  e do 
  \textcolor{blue}
  {\href{https://docs.mapstore.geosolutionsgroup.com/en/v2024.01.02/user-guide/exploring-dashboards/}
  {MapStore}}.

\section{Criação} \label{sec:firstpage}

  Para criar um novo Dashboard siga até a aba "Dados/Conjuntos de Dados" ou
  "Dashboards" no menu superior esquerdo (Figura \ref{fig:criar}.1). No
  catálogo de dados ou de dashboards clique em "Adicionar Recurso/Criar
  dashboard" (Figura \ref{fig:criar}.2).

  \begin{figure}[ht]
    \centering
    \includegraphics[width=\textwidth, keepaspectratio]{img/criarGeostory.pdf}
    \caption{Fluxo para a criação de um novo Geostory.}
    \label{fig:criar}
  \end{figure}

  Uma vez na tela de criação de dashboards, a interface estará dividida em dois
  contêineres. Na esquerda é possível adicionar e modificar os parametros dos
  widgets, enquanto na direita é possível posicionar e visualizar os widgets no
  dashboard. Para criar um novo widget utilize o botão de "+" na parte superior
  esquerda da tela e selecione o tipo de widget que deseja criar.

  \begin{figure}[ht]
    \centering
    \includegraphics[width=\textwidth, keepaspectratio]{img/overviewGeostory.pdf}
    \caption{Interface gráfica de criação do Geostory em destaque o botão de mais que adiciona seções à página.}
    \label{fig:overview}
  \end{figure}

\section{Tipos de Widget}

\subsection{Mapa}

  Widgets de mapa permitem adicionar mapas previamente construidos no Geonode
  ao Dashboard, o mapa pode ser usado para visualização e também para a
  \b{filtragem} dos dados por meio das conexões (ADICIONAR LINK PARA AS
  CONEXOES AQUI).

  Para criar um widget de mapa selecione "Mapa" durante a criação de um novo
  widget e em seguida clique no mapa do Geonode que deseja utilizar, você pode
  usar o campo de texto para buscar os mapas. Em seguida, clique em continuar,
  será possível configurar as camadas visíveis ou adicionar novas camadas ao
  mapa. Também é possível desabilitar a identificação de elementos com o clique
  no mapa. Por fim clique em adicionar widget. Clique em continuar, adicione um
  título e descrição caso desejar, e clique em adicionar widget.

  \begin{figure}[ht]
    \centering
    \includegraphics[width=\textwidth, keepaspectratio]{img/tituloGeostory.pdf}
    \caption{Interface de edição da seção de título do Geostory.}
    \label{fig:titulo}
  \end{figure}

\subsection{Legenda}

  O widget de legenda adiciona uma legenda interativa ao dashboard que deve
  estar obrigatoriamente conectada a um mapa. Além de representar os dados, o
  widget de legenda também permite alternar a visualização das camadas de um
  mapa.

  Para criar um widget de legenda selecione "Legenda" durante a criação de um
  novo widget e em seguida selecione o qual a legenda está conectada, caso o
  dashboard possua mais de um mapa. Clique em continuar, adicione um título e
  descrição caso desejar, e clique em adicionar widget.

  \begin{figure}[ht]
    \centering
    \includegraphics[width=\textwidth, keepaspectratio]{img/midiaGeonode.pdf}
    \caption{Interface gráfica de inserção de mídias do Geonode.}
    \label{fig:midiageonode}
  \end{figure}

  \begin{figure}[ht]
    \centering
    \includegraphics[width=\textwidth, keepaspectratio]{img/midiaExterno.pdf}
    \caption{Interface gráfica de inserção de mídias externas.}
    \label{fig:midiaexterna}
  \end{figure}

\subsection{Gráfico}

  O widget de gráfico adiciona gráficos interativos ao dashboard que podem
  estar ou não correlacionados com os dados demonstrados em mapa.
  
  Para criar um widget de gráfico selecione "Gráfico" durante a criação de um
  novo widget e em seguida selecione o dado que deseja representar, este dado
  deve ser inserido previamente no geonode como um dataset. Em seguida,
  selecione o tipo de gráfico que deseja inserir: barra, pizza ou linha. Na
  tela seguinte é possível editar os aspectos do gráfico: os atributos que irão
  ocupar os eixos x e y, o tipo de operação que será realizado para ingerir os
  dados do eixo y, a cor do gráfico, o uso de legenda e por fim o campo de
  edição avançada que permite alterar a representação do gráfico. Além disso, o
  widget permite o uso de conexões e filtros. Ao finalizar a edição, clique em
  continuar, adicione um título e descrição caso desejar, e clique em adicionar
  widget.

\subsection{Contador}

  O widget de contador adiciona um contador interativo ao dashboard que
  registra uma contagem ou operação matemática que está relacionada ou não aos
  dados do mapa.
  
  Para criar um widget de contador selecione "Contador" durante a criação de um
  novo widget e em seguida selecione o dado que deseja representar, este dado
  deve ser inserido previamente no geonode como um dataset. Na tela seguinte é
  possível editar os aspectos do contador como o atributo a ser representado, a
  operação matemática a ser realizada nos dados deste atributo e a unidade de
  medida do contador. O widget permite o uso de conexões e filtros. Ao
  finalizar a edição, clique em continuar, adicione um título e descrição caso
  desejar, e clique em adicionar widget.

\subsection{Tabela}
  
  O widget de tabela permite a visualização de um dataset na forma de uma
  tabela interativa.
  
  Para criar um widget de table selecione "Tabela" durante a criação de um novo
  widget e em seguida selecione o dataset que deseja representar, este dado
  deve ser inserido previamente no geonode. Em seguida, é possível selecionar
  quais colunas da tabela estarão visiveis no widget. O widget permite o uso de
  conexões e filtros. Por fim, clique em continuar, adicione um título e
  descrição caso desejar, e clique em adicionar widget.

\subsection{Texto}
  O widget de texto permite a inserção de uma caixa de texto ao dashboard.

  Para criar um widget de texto selecione "Texto" durante a criação de um novo
  widget e em seguida insira o texto que deseja representar e opcionalmente um
  título. Também é possível inserir imagens e incorporar mapas e vídeos por
  meio de embeded links. Por fim, clique em adicionar widget.

\section{Conexões}

Os dashboards do geonode permitem a conexão entre os widgets de mapas e tabelas
com os demais widgets, criando assim widgets interativos onde é possível
filtrar os dados que estão sendo representados.

\subsection{Conexão com mapas}

  Ao conectar um mapa a um outro widget, os dados presentes na área que está
  sendo visualizada no mapa serão os dados utilizados pelos demais widgets. De
  forma que, à medida que o usuário altera a visualização do mapa para uma
  região específica os demais widgets irão filtrar os dados para representar
  aquela região.

\subsection{Conexão com tabelas}

  Ao conectar uma tabela a um outro widget, uma linha de filtros será
  adicionada ao topo da tabela. Os filtros aplicados a tabela irão afetar
  também os demais widgets, com exceção dos mapas, que irão mover seu centro
  para representar os dados filtrados na tabela.

\subsection{Criando conexões}

  Durante a inserção de um novo widget, caso haja no dashboard um mapa ou
  tabela no qual o widget pode se conectar um sinal de conexão estará
  disponível. Basta clicar no sinal de conexão e selecionar a qual widget
  deseja conectar (caso apenas um widget conectável esteja disponível, a
  conexão é feita automaticamente). Caso adicione um widget conectável após a
  inserção de outro widget, basta clicar no menu de opções do widget, editar e
  acionar o sinal de conexão.

  \begin{figure}[ht]
    \centering
    \includegraphics[width=\textwidth, keepaspectratio]{img/geocarrousselGeostory.pdf}
    \caption{Fluxo de criação de um Geocarousel.}
    \label{fig:geocarousel}
  \end{figure}


\section{Filtros}
  O construtor permite que o usuário reorganize a ordem das seções do Geostory
  bem como navegue por ele. Utilizando o botão de arrastar no canto superior
  esquerdo de cada seção é possível modificar a ordem dos elementos criados
  anteriormente e utilizando o ícone da lupa no canto direito de cada seção é
  possível mover a visualização até aquele elemento.

  Na parte superior do construtor o usuário pode apagar uma seção selecionada
  através do ícone da lixeira. Uma prévia do Geostory pode ser acessada através
  do ícone de previsão, para retornar ao editor utilize o ícone de edição no
  canto superior esquerdo da página (Figura \ref{fig:previa}).

  \begin{figure}[ht]
    \centering
    \includegraphics[width=\textwidth, keepaspectratio]{img/visualizacaoGeostory.pdf}
    \caption{Entrada e saída do modo de prévia.}
    \label{fig:previa}
  \end{figure}

\section{Configuração}
  Na parte superior do construtor através do ícone de engrenagem é possível
  acessar o menu de configuração. Este permite que o usuário modifique a
  temática geral do Geostory, como cores do plano de fundo, fonte e cores do
  texto e sombreamento das caixas de texto. O usuário também pode adicionar
  elementos na interface gráfica como título da página, logo e atalhos de
  navegação. Cada seção selecionada nas configurações irá criar um atalho de
  navegação no topo da página (Figura \ref{fig:navegacao}).

  \begin{figure}[ht]
    \centering
    \includegraphics[width=\textwidth, keepaspectratio]{img/navegacaoGeostory.pdf}
    \caption{Configuração da barra de navegação de um Geostory.}
    \label{fig:navegacao}
  \end{figure}

\end{document}
